\section*{Introduction}
\addcontentsline{toc}{section}{Introduction}
\label{sec:intro}

\subsection*{Motivation and Objectives}
Exploratory Data Analysis (EDA) is the indispensable first step of any data science pipeline. Before fitting predictive models or designing feature engineering strategies, a practitioner must deeply understand the structure, quality, and statistical characteristics of raw data. Skipping EDA risks building models on corrupted, imbalanced, or poorly understood inputs — leading to misleading evaluation metrics and fragile deployments.

This report documents a systematic EDA conducted across \textbf{four distinct data modalities}, demonstrating how the analytical toolkit must be adapted to the unique properties of each data type:

\begin{enumerate}
    \item \textbf{Tabular data} requires statistical summaries, correlation analysis, and survival/regression-style group comparisons.
    \item \textbf{Text data} demands vocabulary inspection, frequency analysis, stopword handling, and n-gram pattern extraction.
    \item \textbf{Image data} calls for pixel-level statistics, channel distribution analysis, per-class mean images, and data quality audits.
    \item \textbf{Multimodal data} involves exploring both the visual and linguistic modalities in tandem, checking alignment and coverage.
\end{enumerate}

\subsection*{Datasets at a Glance}

\begin{importantbox}
\textbf{Four Modalities — Four Benchmark Datasets:}
\begin{itemize}
    \item \textbf{Tabular} \textrightarrow\ \textbf{Titanic Passenger Dataset:} 891 passenger records, 12 features, binary classification target (\texttt{Survived}).
    \item \textbf{Text} \textrightarrow\ \textbf{SMS Spam Collection:} 5{,}572 labeled SMS messages, binary classification (\texttt{ham} / \texttt{spam}).
    \item \textbf{Image} \textrightarrow\ \textbf{CIFAR-10:} 50{,}000 color images ($32{\times}32$ px), 10 perfectly balanced object categories.
    \item \textbf{Multimodal} \textrightarrow\ \textbf{Flickr8k:} 8{,}092 real-world photographs each paired with 5 human-annotated English captions (40{,}460 captions total).
\end{itemize}
\end{importantbox}

\subsection*{Analysis Framework}

For each dataset, we apply a consistent four-stage framework:

\begin{enumerate}
    \item \textbf{Structure Inspection} — Dataset shape, column types, memory footprint, first-row preview.
    \item \textbf{Quality Audit} — Missing values, duplicate records, out-of-range pixel values, label consistency.
    \item \textbf{Distribution Analysis} — Univariate histograms, class balance checks, per-channel pixel statistics.
    \item \textbf{Relationship Analysis} — Correlation matrices, survival group comparisons, class-conditional word frequencies, inter-annotator agreement (for multimodal).
\end{enumerate}

All analysis was performed in Python using \texttt{pandas}, \texttt{numpy}, \texttt{matplotlib}, \texttt{seaborn}, \texttt{wordcloud}, \texttt{tensorflow.keras}, and \texttt{datasets} (HuggingFace), and is fully reproducible via the accompanying Jupyter notebook.
